\documentclass[9pt,landscape]{article}
\usepackage[letterpaper,margin=0.42in]{geometry}
\usepackage{amsmath}
\usepackage{booktabs}
\usepackage{graphicx}
\usepackage{xcolor}
\usepackage{microtype}
\usepackage{enumitem}
\usepackage[T1]{fontenc}
\usepackage{lmodern}

\definecolor{pairwise}{HTML}{B42318}
\definecolor{unary}{HTML}{374151}
\definecolor{ink}{HTML}{111827}
\definecolor{muted}{HTML}{4B5563}
\definecolor{panel}{HTML}{F3F4F6}
\definecolor{blue}{HTML}{355F8A}

\newcommand{\StepSevenGainRows}{%
GSM8K & +0.418 & +0.444 & +0.436 \\
MATH500 & +0.404 & +0.439 & +0.434 \\
}
\newcommand{\ExactMatchCount}{768}
\newcommand{\ExactMatchTotal}{768}


\pagestyle{empty}
\setlength{\parindent}{0pt}
\setlength{\parskip}{2.2pt}
\setlist[itemize]{leftmargin=10pt,itemsep=1.2pt,topsep=1.5pt}
\newcommand{\heading}[1]{%
  \vspace{1pt}{\color{ink}\bfseries\large #1}\vspace{1pt}\par}
\newcommand{\smallheading}[1]{%
  {\color{ink}\bfseries #1}\hspace{2pt}}
\newcommand{\pair}{\textcolor{pairwise}{Pairwise-K16}}
\newcommand{\unaryname}{\textcolor{unary}{Unary-K16}}

\begin{document}
\color{ink}

\begin{minipage}[t]{0.73\textwidth}
  {\fontsize{20}{22}\selectfont\bfseries Conditional Draft Trees with DFlash2}\\[-1pt]
  {\fontsize{10.5}{12}\selectfont\color{muted}
  Using predecessor-conditioned scores to allocate speculative verification trees}
\end{minipage}
\hfill
\begin{minipage}[t]{0.25\textwidth}
  \raggedleft\footnotesize
  Research exploration\\
  Frozen Qwen3-4B + official Qwen3.8-27B evidence
\end{minipage}

\vspace{2pt}\hrule\vspace{5pt}

\begin{minipage}[t]{0.485\textwidth}
\heading{Problem and idea}
\colorbox{panel}{%
  \parbox{\dimexpr\linewidth-2\fboxsep\relax}{%
  \smallheading{Problem.}
  DDTree allocates a fixed \(B\)-node verification tree using independent
  positional marginals. The score of \emph{York} at depth 2 is unchanged
  after \emph{New} or \emph{Los}.

  \smallheading{Idea.}
  DFlash2 already exposes predecessor-conditioned selector scores.
  Convert them into bounded, mass-preserving transitions and reuse the same
  DDTree best-first allocator and target verifier.}}

\vspace{3pt}
\begin{minipage}[t]{0.47\linewidth}
  \centering
  \textbf{\color{unary}Unary allocation}\\[-2pt]
  \[
    Q_U(y_{1:d})=\prod_{t=1}^{d}q_t(y_t)
  \]
  \scriptsize Independent positional marginals
\end{minipage}
\hfill
\begin{minipage}[t]{0.51\linewidth}
  \centering
  \textbf{\color{pairwise}Conditional allocation}\\[-2pt]
  \[
    Q_P(y_{1:d}\!\mid y_0)
    =\prod_{t=1}^{d}m_t\,r_t(y_t\!\mid y_{t-1})
  \]
  \scriptsize \(y_0\): verified anchor; \(m_t\): retained top-16 unary mass
\end{minipage}

\vspace{2pt}
\begingroup\small
\[
\log m_t=\operatorname{LSE}_{b\in C_t}U_t(b)
-\operatorname{LSE}_{v\in\mathcal V}U_t(v),
\qquad
r_t(b\mid a)=\operatorname{softmax}_{b\in C_t}
\!\left[U_t(b)+c_t(a,b)\right].
\]
\endgroup
\footnotesize
\(m_t\) preserves how much full-vocabulary mass lies in \(C_t\);
\(r_t\) redistributes it using DFlash2's learned selector.
Every extension lies in \([0,1]\), so descendants cannot outrank parents and
best-first construction remains prefix-closed.

\heading{Why conditioning can spend \(B\) better}
\begin{center}
\begin{tabular}{@{}rcl@{\qquad}rcl@{}}
\textbf{New} & \(\rightarrow\) &
  \textcolor{pairwise}{\textbf{York}} & 
\textbf{Los} & \(\rightarrow\) &
  \textcolor{pairwise}{\textbf{Angeles}} \\
 & \(\searrow\) & Jersey &
 & \(\searrow\) & \textcolor{muted}{York}
\end{tabular}
\end{center}
\vspace{-3pt}
\scriptsize
Unary knows that \emph{York} is individually likely. Conditional allocation
can prefer coherent paths (\emph{New York}, \emph{Los Angeles}) over
inconsistent combinations. The contribution changes \textbf{tree allocation},
not target verification.

\heading{Falsification test: candidate breadth}
\vspace{-3pt}
\begin{center}
\includegraphics[width=\linewidth]{figures/wider_unary_falsification.pdf}
\end{center}
\vspace{-8pt}
\scriptsize
\textbf{Qwen3-4B, Step 6.3.} Wider unary support improves target-path
representability but barely changes fixed-\(B\) acceptance; \pair{} remains
above Unary-K16/K32/K64. \(K\) is candidate support; \(B\) is verified nodes.
\end{minipage}
\hfill
\begin{minipage}[t]{0.485\textwidth}
\heading{Transfer across checkpoint and model scale}
\vspace{-4pt}
\begin{center}
\includegraphics[width=\linewidth]{figures/model_scale_transfer.pdf}
\end{center}
\vspace{-8pt}
\scriptsize
Points show \pair{} minus \unaryname{} mean matched draft tokens;
bars are paired prompt-bootstrap 95\% CIs.

\vspace{4pt}
\begin{minipage}[t]{0.54\linewidth}
\smallheading{Official Qwen3.8-27B DFlash2}
\vspace{2pt}
\centering
\begin{tabular}{@{}lrrr@{}}
\toprule
Pairwise gain & \(B{=}16\) & \(B{=}32\) & \(B{=}64\) \\
\midrule
\StepSevenGainRows
\bottomrule
\end{tabular}

\vspace{3pt}
\raggedright
\textbf{All paired 95\% CIs exclude zero.}\\
\(\ExactMatchCount/\ExactMatchTotal\) Unary/Pairwise outputs exactly matched
sequential target decoding.
\end{minipage}
\hfill
\begin{minipage}[t]{0.42\linewidth}
\smallheading{Cross-domain 4B evidence}
\begin{itemize}
  \item Same \(K=16\), same \(B\): Pairwise improves matched tokens on
    GSM8K, MATH500, HumanEval, and MT-Bench.
  \item The 4B prototype converts higher acceptance into positive throughput
    gains at every tested budget/domain.
\end{itemize}
\end{minipage}

\heading{What the evidence supports}
\begin{itemize}
  \item \textbf{Conditional allocation is consistently better under the
    tested fixed-node budgets.}
    Wider unary support to \(K=64\) does not recover the gain.
  \item \textbf{The effect transfers} from an experimental Qwen3-4B DFlash2
    checkpoint to the official Qwen3.8-27B DFlash2 pair.
  \item \textbf{System claim is architecture-specific.}
    The 27B correctness-first recurrent verifier cannot parallelize branches;
    Pairwise and Unary throughput are statistically tied there.
\end{itemize}

\heading{Mathematical boundary}
\footnotesize
For the explicit conditional surrogate (augmented with residual mass outside
top-16), DDTree's tail-sum objective remains additive and best-first returns
the highest-mass \(B\) prefixes on the retained lattice.
\textbf{We do not claim} optimality under the target distribution, unrestricted
vocabulary optimality, or that arbitrary conditional scores improve ranking.

\vfill
\hrule\vspace{3pt}
\scriptsize\color{muted}
\textbf{Attribution.}
DDTree: verification-tree objective, best-first construction, verifier.
DFlash2: drafter, top-16 lattice, predecessor-conditioned selector.
This extension: mass-preserving conditional scorer, integration, controlled
falsification and transfer studies.\\
\textbf{Frozen sources.}
Step 6.2, Step 6.3, and Step 7 committed analysis tables; figures are generated
directly from their CSV/JSON artifacts.
\end{minipage}

\end{document}
